\section{Опис побудови коду індивідуального завдання. Варіант №7}

\subsection{Об'єктно-орієнтована декомпозиція та архітектура колекцій}

Згідно з індивідуальним Варіантом №7, було розроблено об'єктно-орієнтовану програму 
для комплексного обліку та аналізу фільмів у кінокаталозі. На відміну від першої 
частини роботи, де операції виконувалися безпосередньо у методі \texttt{main()}, 
тут всю бізнес-логіку обробки контейнерів повністю інкапсульовано всередині окремого 
сервісного класу-менеджера.

Архітектура індивідуального завдання побудована на взаємодії таких компонентів:
\begin{enumerate}
    \item \textbf{Клас сутності \texttt{Movie.java}:} 
    Описує модель фільму, повністю приховуючи (інкапсулюючи) внутрішні поля назви 
    (\texttt{title}), жанру (\texttt{genre}), року випуску (\texttt{year}) та 
    рейтингу (\texttt{rating}) за допомогою модифікатора \texttt{private}. Доступ 
    до атрибутів надається через публічні методи-геттери.

    \item \textbf{Клас логіки \texttt{MovieManager.java}:} 
    Виступає в ролі локальної бази даних предметої області. Всередині нього 
    оголошено дві узагальнені структури даних: лінійний список \texttt{List<Movie>} 
    для збереження хронології додавання елементів та хеш-множину \texttt{Set<String>} 
    для накопичення жанрів.
\end{enumerate}

\subsection{Аналіз роботи узагальнених структур даних (\textit{Generics})}

Впровадження колекцій \textit{Java Collections Framework} дозволило вивести 
надійність та гнучкість коду індивідуального завдання на промисловий рівень за 
рахунок таких рішень:
\begin{itemize}
    \item \textbf{Автоматична фільтрація унікальних атрибутів:} 
    При додаванні нового об'єкта через метод \texttt{addMovie(Movie movie)}, 
    менеджер не просто вставляє посилання у список, а паралельно виконує 
    інструкцію \texttt{uniqueGenres.add(movie.getGenre())}. Це дозволяє динамічно 
    формувати чисту множину унікальних жанрів кінематографа «на льоту», повністю 
    позбавляючи розробника від написання громіздких вкладених циклів порівняння рядків.

    \item \textbf{Типобезпечний пошук лідера рейтингу:} 
    Завдяки використанню суворих узагальнень (\texttt{List<Movie>}), компілятор 
    Java гарантує, що у списку містяться виключно об'єкти фільмів. Це дозволило 
    реалізувати швидкий лінійний алгоритм пошуку максимального рейтингу за допомогою 
    методу \texttt{get(0)} та прямого виклику геттера \texttt{movie.getRating()} без 
    використання ризикованих операцій приведення типів та оператора \texttt{instanceof}, 
    які застосовувалися у минулих лабораторних роботах.
\end{itemize}
